\documentclass{article}
\usepackage{geometry}
\usepackage{hyperref}
\usepackage{amsfonts}
\usepackage{amsmath}
\usepackage{graphicx}
\usepackage{/home/user1/code/tex/sty}
\setlength{\parindent}{0pt}
\geometry{top=1in,
  bottom=1in,
  left=1in,
  right=1in,
}
\newcommand{\pic}[2]{\begin{center}\includegraphics[height=#1]{#2}\end{center}}
\newcommand{\abs}[1]{|#1|}
\usepackage{graphicx}
\n{lorentz}{(1-(\fr{v}{c})^2)^{-1/2}}
\n{sxx}{\sigma_{xx}}
\n{syy}{\sigma_{yy}}
\n{sxy}{\sigma_{xy}}
\n{syx}{\sigma_{yx}}

\begin{document}
\begin{center}
  {\Large{Homework 3}}\\~\\
  \begin{tabular}{l l}
    Student & Agustin Romero\\
    Instructor & Sri Raghu\\
    Assistants & Josephine Yu, Pavel Nosov\\
    Course & Classical Electrodynamics\\
    Institution & Stanford University\\
    Due Date & 17 February 2023\\
  \end{tabular}
\end{center}
Collaborators:
\beg{enumerate}{
  \item Charles Yang
  \item Andrew Sullivan
}
\section{Question 1.a.}
\pic{5cm}{1.jpg}
Ref. Griffiths p75, p103. In the plate frame, the surface charge is not moving. The magnetic field everywhere is 0.
\e{\boxed{\vec{B}=0}}
Surfaces with conductivity $\sigma$ produce the field
\e{\vec{E}=\fr{\sigma}{\epz}\hat{n}}
The electric field produced by the $+\sigma$ plate is
\e{E=\fr{1}{2\epsilon_0}\sigma}
pointing away itself. The $-\sigma$ plate produces the field
\e{E=\fr{1}{2\epsilon_0}\sigma}
pointing towards itself. Between the plates, the two fields add,
\e{\boxed{\vec{E}_\rom{inside}=0\hat{x}-\fr{\sigma}{\epsilon_0}\hat{y}+0\hat{z}}}
Outside the plates, the two fields cancel.
\e{\boxed{\vec{E}_\rom{outside}=0}}

\section{Question 1.b.i.}
Ref. Griffiths p553. Let the rocket frame be $O'$ in the $\hat{x}$ direction. Gauss' equation
\e{\int \vec{E}d\vec{a}=\fr{1}{\epz}Q_{\rom{inc}}}
is valid for charges with constant velocity, so the analysis from 1.a. still applies. $\vec{E'}$ is in the \fbox{$-\hat{y}$} direction with no $\hat{x}$ nor $\hat{z}$ components.
Next, consider the $\vec{B'}$ field. Ampere's equation is
\e{\nabla\times \vec{B'}=\mu_0 \vec{j'}}
$\vec{j'}$ is in the $-\hat{x}$ direction, so ${\vec{B'}}$ is in the \fbox{+$\hat{z}$} direction with no $\hat{x}$ nor $\hat{y}$ components.

\section{Question 1.b.ii.}
The question didn't state whether to use Lorentz or Galilean transformations. Choose Lorentz. Solve for $B_z$.
\e{B_z=\muz \sigma v}
Solve for $B_z'$.
\e{B_z'&=\lorentz (B_z-\fr{v}{c^2}E_y)\\
&=\lorentz(\muz \sigma v + \fr{v}{c^2} \fr{\sigma}{\epz})}
The units are consistent.
\e{[\muz][\sigma][v]=T}
\e{[v][c]^{-2}[\sigma][\epz]^{-1}=T}
Solve for $E_y'$.
\e{E_y'&=\lorentz (E_y - v B_z)\\
&=\lorentz (-\fr{\sigma}{\epz}-v^2 \muz \sigma) \hat{y}}
$\vec{E'}$ is
\e{\boxed{\vec{E'}=0\hat{x}-\lorentz (\fr{\sigma}{\epz} + v^2 \muz \sigma)\hat{y}+0\hat{z}}}
$\vec{B'}$ is
\e{\boxed{\vec{B'}=0\hat{x}+0\hat{y}+\lorentz(\muz \sigma v + \fr{v \sigma}{c^2 \epz})\hat{z}}}

\section{Question 1.c.i.}
Let the rocket frame be $O'$ in the $\hat{y}$ direction. Gauss' equation still applies. $\vec{E'}$ is in the \fbox{$-\hat{y}$} direction, with no $\hat{z}$ nor $\hat{x}$ components. Next, consider the $\vec{B'}$ field. The rocket observes one plate moving towards it with velocity $-\vec{v}\hat{y}$, and the other plate moving away with velocity $+\vec{v}\hat{y}$. The rocket observes that neither plate has a surface current $\vec{j}$, therefore Ampere's equation is
\e{\nabla \times \vec{B'} = 0}
and \fbox{$\vec{B'}=0$}.

\section{Question 1.c.ii.}
In the boosted frame, the $\vec{E}$ field is nonzero in the direction towards the plates, $\hat{y}$, but the $\vec{B}$ fields is zero.
\e{\boxed{\vec{E'}=0\hat{x}+\lorentz \fr{\sigma}{\epz}\hat{y}+0\hat{z}}}
\e{\boxed{\vec{B'}=0\hat{x}+0\hat{y}+0\hat{z}}}

\section{Question 1.d.}The boost equations are
\e{E'_{||}=E_{||}}
\e{B'_{||}=B_{||}}
\e{E'_\perp=\lorentz (E_\perp+\vec{v}\times\vec{B})}
\e{B'_\perp=\lorentz (B_\perp-\fr{\vec{v}}{c^2}\times\vec{E})}
Let $\vec{v}$ be in the $\hat{x}$ direction. The component parallel to the boost, $\hat{x}$, is not transformed by the boost. For $E'_x$ to be 0, $E_x$ must be 0. The $\hat{y}$ and $\hat{z}$ components of $\vec{E}$ and $\vec{B}$ mix, and must be exactly negative of the resulting boosted fields $\vec{E'}$ and $\vec{B'}$. $\vec{v}\times \vec{B}$ is
\e{\vec{v}\times \vec{B}&=(v_yB_z-B_yv_z)\hat{x}-(v_xB_z-B_xv_z)\hat{y}+(v_xB_y-B_xv_y)\hat{z}\\
&=-v_xB_z\hat{y}+v_xB_y\hat{z}}
The condition is
\e{\label{noqwid}\boxed{\vec{E}=0\hat{x}+v_x B_z\hat{y}-v_x B_y\hat{z}}}

\section{Question 1.e.}The logic is identical to 1.d..
\e{\vec{v}\times \vec{E}&=(v_yE_z-E_yv_z)\hat{x}-(v_xE_z-E_xv_z)\hat{y}+(v_xE_y-E_xv_y)\hat{z}\\
&=-v_xE_z\hat{y}+v_xE_y\hat{z}}
The condition is
\e{\label{aoduna}\boxed{\vec{B}=0\hat{x}-\fr{v_x}{c^2}E_z\hat{y}+\fr{v_x}{c^2}E_y\hat{z}}}

\section{Question 1.f.}
\pic{5cm}{2}
The Hall experiment can be summarized as:
\beg{enumerate}{
\item Setup two magnets, a line of current, and a dielectric.
\item Send current through the dielectric.
\item Observe the current is temporarily perturbed, and a voltage V appears in the $\hat{z}$ direction.
\item Observe the current returns to its unperturbed path, and the voltage V is constant.
}
First, think experimentally. In this setup, $\vec{E}$ points in the $\hat{x}$ direction, $\vec{B}$ in the $\hat{y}$ direction. The Lorentz force implies that the current in the dielectric is forced in the $+\hat{z}$ direction, generating a voltage V, until the force $\vec{F}=q\vec{E}$ becomes equal and opposite to the force $\vec{v}\times\vec{B}$, thereby returning the current to its unperturbed path. Because the fields are constrained along a plane, the condition that the path is unperturbed is
\e{\label{o12n39812}E_y=v_x B_z}
which is exactly the result found in the previous parts. Next, think deductively. The result \eref{o12n39812} was derived by claiming there should exist an inertial frame where the $\vec{E'}$ does not exist. The equation
\e{\vec{E}\cdot\vec{B}=0}
holds, and so the results of 1.d. and 1.e should apply. Observe that 1.d. predicts that \eref{noqwid} must be satisfied for there to exist a frame where $\vec{E}$ is zero, and 1.e. predicts that \eref{aoduna} must be satisfied for there to exist a frame where $\vec{B}$ is zero. Observe that the Hall experiment can satisfy either of these equations, where $v_x$ is the velocity of the electrons. The conclusion is that \fbox{all frames agree on the experimental predictions}, where in one frame the current is unaffected by the $\vec{E}$ and $\vec{B}$ fields, and in another frame the fields affect the current equally in magnitude but opposite in direction. One claim of special relativity is that experimental predictions are the same in all reference frames, and that is true here.

\section{Question 2.a.}
First of all,
\e{[E]=[E]+[1/c][B][v]=[E]+[B]}
So there seems to be an issue with the equations. For the sake of sanity modify the equations so that $\vec{j}$ will have the correct units. Solve for $\vec{B'}$.
\e{\vec{B}=\vec{B'}-\fr{1}{c^2}\vec{E'}\times \vec{v}}
\e{\vec{B}=\vec{B'}}
\e{\vec{E}&=\vec{E'}+\vec{B'}\times \vec{v}\\
&=((B'_yv_z-B_z'v_y)\hat{x}-(B_x'v_z-B_z'v_x)\hat{y}+(B_x'v_y-B_y'v_x)\hat{z})
\intertext{$B_x'=B_x=0$ and $B_y'=B_y=0$.}
&=\fr{1}{c}(-B_z'v_y\hat{x}+B_z'v_x\hat{y}+0\hat{z})
\intertext{$E_x=0$.}
&=\label{iundia}(0\hat{x}+B_z'v_x\hat{y}+0\hat{z})}
\eref{iundia} determines $v_y$.
\e{0=-B_z'v_y}
Because $B_z'\neq 0$, $v_y=0$. Because $E_y\neq 0$,
\e{v_x=\fr{E_y}{B_z'}}
The magnetic field in the boosted frame is
\e{\boxed{\vec{B'}=0\hat{x}+0\hat{y}+B_z'\hat{z}=0\hat{x}+0\hat{y}+\fr{E_y}{v_x}\hat{z}}}
To find $v_z$, observe that the electrons are in the xy plane, so $v_z=0$. The boost velocity vector is
\e{\boxed{\vec{v}=\fr{E_y}{B_z}\hat{x}+0\hat{y}+0\hat{z}}}

\section{Question 2.b.}
Using the result for $\vec{v}$ $\vec{j}$ is
\e{\boxed{\vec{j}=ne\vec{v}=ne(\fr{E_y}{B_z}\hat{x}+0\hat{y}+0\hat{z})}}
Because $\vec{E}$ is in the $+\hat{y}$ direction, $\vec{j}$ is \fbox{perpendicular} to $\vec{E}$.

\section{Question 2.c.}
Previously,
\e{\vec{E}=0\hat{x}+B_zv_x\hat{y}+0\hat{z}}
and
\e{\vec{j}=ne\fr{E_y}{B_z}\hat{x}+0\hat{y}+0\hat{z}}
so
\e{j_x=\sxx E_x + \sxy E_y = \sxy E_y = ne\fr{E_y}{B_z}}
\e{\sxy=\fr{ne}{B_z}}
\e{\syx=\fr{-ne}{B_z}}
\e{j_y=\syx E_x + \syy E_y = \syy E_y = 0}
\e{j_y=0\quad E_y\neq 0 \Rightarrow \syy=0}
\e{\boxed{\sxx=0\quad \syy=0}}

\section{Question 2.d.}
Redo the previous parts using a Lorentz transformation. The only assumptions are $\vec{E}$, $\vec{B}$, and $\vec{E'}=0$.
\e{\vec{E}=0\hat{x}+E_y\hat{y}+0\hat{z}}
\e{\vec{B}=0\hat{x}+0\hat{y}+B_z\hat{z}}
The undetermined components are $E_y$ and $E_z$.
\e{E_y'=\lorentz (E_y+vB_z)}
\e{E_y=-\fr{v}{c}B_z}
\e{E_z'=\lorentz (E_z+vB_y)=0}
\e{B_x'=B_x=0}
\e{B_y'=\lorentz (B_y-\fr{v}{c^2}E_z)=0}
\e{B_y=\fr{v}{c^2}E_z}
\e{B_z'&= \lorentz (B_z+\fr{v}{c^2}E_y)\\
&=\lorentz (B_z-\fr{v^2}{c^2}B_z)}
\e{\vec{j}=ne\vec{v}=ne\fr{E_y}{B_z}\hat{x}+0\hat{y}+0\hat{z}}
\e{j_x&=\sxx E_x + \sxy E_y\\
&=\sxy E_y\\
&= -\fr{ne}{B_z}E_y}
\e{\sxy = - \fr{ne}{B_z}}
\e{j_y=\syx E_x + \syy E_y = \syy E_y = 0}
\e{E_y=0\Rightarrow \syy=0}
\e{\sxx=0}
The resistivity tensor is
\e{\boxed{\sxx = 0 \quad \syy = 0 \quad \sxy = -\fr{ne}{B_z} \quad \syx = \fr{ne}{B_z}}}
The $\vec{B'}$ field is
\e{\boxed{\vec{B'}=0\hat{x}+0\hat{y}+\lorentz (B_z-\fr{v^2}{c^2}B_z) \hat{z}}}
The $\vec{B'}$ field and resistivity tensor $\fbox{both change}$.

\section{Question 3.a.}
\n{deltdtphi}{\delta \dtphi}
\n{deltdxphi}{\delta d_x \phi}
\n{deltphi}{\delta \phi}
\n{dtphi}{d_t \phi}
\n{dxphi}{d_x \phi}
\n{lag}{\mathcal{L}}
\n{oocs}{\fr{1}{c^2}}
\n{padtphi}{\fr{\partial}{\partial \dtphi}}
\n{padt}{\fr{\partial}{\partial t}}
\n{padxphi}{\fr{\partial}{\partial d_x \phi}}
\n{paphi}{\fr{\partial}{\partial \phi}}
\n{patsq}{\fr{\partial^2}{\partial t^2}}
\n{paxsq}{\fr{\partial^2}{\partial x^2}}
\n{pax}{\fr{\partial}{\partial x}}
\n{pat}{\fr{\partial}{\partial t}}
\n{phio}{\phi_1}
\n{phit}{\phi_2}
\n{dt}{\fr{d}{dt}}
\n{dx}{\fr{d}{dx}}
\n{dtsq}{d_t^2}
\n{dxsq}{d_x^2}
\n{notsurface}{\int (\deltphi \paphi \lag - \deltphi \dt \padtphi \lag - \deltphi \dx \padxphi \lag) dx dt}
\n{surface}{\int \padtphi (\padtphi \lag)_{t_1}^{t_2}dx + \int \deltdxphi (\padxphi \lag)_{t_1}^{t_2} dx}
\n{eom}{\paphi \lag - \dt \padtphi \lag - \dx \padxphi \lag}
On notation, $d_x$ is the total derivative, $\deltphi$ is the perturbation of $\phi$, and $\deltdxphi$ is the perturbation of $\dx \phi$.
\e{S_1=\int \lag(\phi, \dtphi, \dxphi) dx dt}
Perturb.
\e{S_2&=\int \lag(\phi + \deltphi, \dtphi +\deltdtphi, \dxphi + \deltdxphi) dx dt\\
&=S_1+\int (\deltphi \paphi \lag + \deltdtphi \padtphi \lag + \deltdxphi \padxphi \lag) dx dt\\
\intertext{Integrate by parts.
\e{\int \deltdtphi \padtphi \lag dx dt=\int \deltphi (\padtphi \lag)_{t_1}^{t_2} dx - \int \deltphi \dt \pat \lag dx dt}
\e{\int \deltdxphi \padxphi \lag dx dt=\int \deltphi (\padxphi \lag)_{t_1}^{t_2} dx - \int \deltphi \dx \pax \lag dx dt}}
&=S_1 + \notsurface + \surface
\intertext{Set the ``surface terms'' to 0.
\e{\surface=0}}
&=S_1 + \notsurface\\
\intertext{Factor out $\deltphi$.}
&=S_1 + \int \deltphi (\eom) dx dt}
The perturbed action gives the unperturbed action, plus a differential equation in $\lag$. It is the Equation of Motion.
\e{\boxed{\eom=0}}
Next, substitute $\lag$ into the EOM to find the EOM in terms of $\phi$.
\e{\lag = \fr{1}{2c^2}(\dtphi)^2-\fr{1}{2}(\dxphi)^2-\fr{1}{2}m^2\phi^2}
\e{\boxed{m^2\phi-\fr{1}{c^2}\dtsq\phi+\dxsq\phi=0}}
Next, assume $\phio$ and $\phit$ solve the EOM, and test if $\phio+\phit$ also solves the EOM.
\e{\phi=\phio+\phit}
\e{-m^2\phio-m^2\phit-\fr{1}{c^2}\dtsq(\phio+\phit)+\dxsq(\phio+\phit)=0}
\e{(-m^2\phio-\fr{1}{c^2}\dtsq\phio+\dxsq\phio)+(-m^2\phit-\fr{1}{c^2}\dtsq\phit+\dxsq\phit)=0}
The linear combination splits into two EOMs, each of which are solved.
\e{\boxed{0+0=0}}
In regards to the minutia, the integration by parts enacts a total derivative in $x$ and $t$, not a partial derivative, which gives the $\dxsq \phi$ and $\dtsq \phi$ terms in the EOM, instead of something like $\pax \dx \phi$ which is incorrect.

\section{Question 3.b.}
Substitute $\lag$ into the EOM to find the EOM in terms of $\phi$.
\e{\lag={\fr{1}{2c^2}\dtphi^2-\fr{1}{2}(\dxphi)^2 - \fr{1}{2} m^2 \phi^2 - \fr{\lambda}{4} \phi^4}}
\e{\label{91283n}\boxed{-m^2\phi-\fr{1}{c^2}\dtsq \phi+\dxsq \phi - \lambda \phi^3=0}}
Next, assume $\phio$ and $\phit$ solve the EOM, and test if $\phio+\phit$ also solves the EOM.
\e{\phi=\phio+\phit}
\e{0&=-m^2\phio-m^2\phit-\oocs \dtsq \phio - \oocs \dtsq \phit + \dxsq \phio + \dxsq \phit - \lambda(\phio + \phit)^3\\
&=(-m^2\phio-\oocs \dtsq \phio+ \dxsq \phio)+(-m^2\phit- \oocs \dtsq \phit+ \dxsq \phit)-\lambda(\phio^3+\phit^3+3\phio^2\phit+3\phit^2\phio)\\
&=-3\lambda(\phio^2\phit+\phit^2\phio)}
\e{\label{ap9w8rhq}\boxed{-3\lambda(\phio^2\phit+\phit^2\phio)=0}}
In order for the equation of motion \eref{91283n} to be satisfied, \eref{ap9w8rhq} must also be satisfied, which means the field $\phi=\phio+\phit$ is \fbox{not necessarily a solution} to the EOM. As a comment, these could be called ``secondary'' equations of motion, and it would be interesting to find out if they are necessary to correctly predict dynamics, especially for theories such as $\lambda \phi^4$.

\end{document}